\thispagestyle{empty}

\begin{minipage}{0.3\textwidth}
  \includegraphics[scale=0.35]{logounal.png}
\end{minipage}%
\hfill
\begin{minipage}{0.65\textwidth}
  \begin{center}
    \scshape
    \Large \textsc{Universidad Nacional de Colombia} \\
    \textcolor{white}{\tiny.} \Large \textsc{Departamento de Matemáticas} \\
    \textcolor{white}{\tiny.} \large \textsc{Análisis Funcional} \\
    \textcolor{white}{\tiny.} \large \textsf{Taller 1: Espacios vectoriales normados} \normalsize (I-2025)
  \end{center}
\end{minipage}

\vspace{0.3cm}
\normalfont

\textbf{Profesor:} Oscar Guillermo Riaño Castañeda\\
\textbf{Integrantes:} Andrés David Cadena Simons \hspace{2.8cm}  Jairo Sebastián Niño Castro\hspace{2.8cm}
Iván Felipe Salamanca Medina \hspace{5.05cm}\textbf{Fecha:} 02 de Julio del 2025\\
\vspace{0.25cm}\\

\section{Espacios de Hilbert}

\textbf{Ejercicio 13.}

\begin{enumerate}
    \item[(I)] Muestre que los siguientes conjuntos $M$ son subespacios cerrados no vacíos de $L^2((-1,1))$ y determine explícitamente la proyección $P_M$ en cada caso
    \begin{enumerate}
        \item[(a)] $M=\{f \in L^2((-1,1)):f(x)=f(-x) \text{ para casi todo } x \in (-1,1)\}$.
        \item[(b)] $\displaystyle M=\left\{f \in L^2((-1,1)):\int_{-1}^1f(x)\, dx=0\right\}$.
        \item[(c)] $M=\{f \in L^2((-1,1)):f(x)=0 \text{ para casi todo } x\in (-1,0)\}$.
    \end{enumerate}
    \item[(II)] Sea $\Omega\subset \mathbb{R}^n$ un abierto acotado. Considere
    \begin{align*}
        K=\left\{f\in L^2(\Omega):\int_{\Omega}f(x)\, dx\geq 1\right\}.
    \end{align*}
    \begin{enumerate}
        \item[(a)] Demuestre que $K$ es un conjunto cerrado convexo en $L^2(\Omega)$.
        \item[(b)] Determine la proyección sobre $K$, es decir, el operador $P_K$.
    \end{enumerate}
    \end{enumerate}

    \begin{proof}
    \begin{enumerate}
        \item[(I)] 
        \begin{enumerate}
            \item[(a)] Sea $(f_n)$ una sucesión en $M$ tal que $f_n\to f\in L^2((-1,1))$. Para todo $\varepsilon>0$ existe $N_\varepsilon\in \mathbb{Z}^+$ tal que si $n\geq N_\varepsilon$, entonces
            \begin{align*}
                \norm{f_n-f}_{L^2((-1,1))}=\left(\int_{-1}^1 |f_n(x)-f(x)|\, dx\right)^{1/2}<\dfrac{\varepsilon}{2},
            \end{align*}
            Para $x \in (-1,1)$, sea $g(x)=f(-x)$ y veamos que $\norm{g-f}_{L^2((-1,1))}=0$. Sea $n\geq N_\varepsilon$, tenemos
            \begin{align*}
                \norm{g-f}_{L^2((-1,1))}\leq \norm{g-f_n}_{L^2((-1,1))}+\norm{f_n-f}_{L^2((-1,1))}.
            \end{align*}
            Como $n\geq N_\varepsilon$, $\norm{f_n-f}_{L^2((-1,1))}<\dfrac{\varepsilon}{2}$. Ahora, usando que $f_n\in M$ y la sustitución $y=-x$, tenemos
        \begin{align*}
            \norm{g-f_n}_{L^2((-1,1))}^2&=\int_{-1}^1|g(x)-f_n(x)|^2\, dx\\
            &=\int_{-1}^1|f(-x)-f_n(x)|^2\, dx\\
            &=\int_{-1}^1|f(y)-f_n(y)|^2\, dy\\
            &=\norm{f-f_n}_{L^2((-1,1))}^2,
        \end{align*}
        de esta manera $\norm{g-f_n}_{L^2((-1,1))}=\norm{f-f_n}_{L^2((-1,1))}<\dfrac{\varepsilon}{2}$, de esta manera
        \begin{align*}
            \norm{g-f}_{L^2((-1,1))}<\dfrac{\varepsilon}{2}+\dfrac{\varepsilon}{2}=\varepsilon.
        \end{align*}
        Como $\varepsilon>0$ es arbitrario, tenemos que $\norm{g-f}_{L^2((-1,1))}=0$, lo que quiere decir que $f(-x)=g(x)=f(x)$ para casi todo $x \in (-1,1)$, es decir, $f \in M$ y de esta manera $M$ es cerrado.
        
        Veamos ahora que $M$ es un subespacio. Sean $f,g \in M$ y $\alpha\in \mathbb{R}$. Sean 
        \begin{align*}
            E_f&:=\{x \in (-1,1):f(x)\neq f(-x)\}\\
            E_g&:=\{x \in (-1,1):g(x)\neq g(-x)\},
        \end{align*}
        como $f,g \in M$, 
        \begin{align*}
            0\leq \mu(E_f\cup E_g)\leq \mu(E_f)+\mu(E_g)=0+0=0.
        \end{align*}
        es decir, $\mu(E_f\cup E_g)=0$, además, para todo $x \in (-1,1)\setminus (E_f\cup E_g)$ se tiene que $f(x)=f(-x)$ y $g(x)=g(-x)$, de manera que 
        \begin{align*}
            f(x)+g(x)=f(-x)+g(-x) \text{ para todo } x \in (-1,1)\setminus (E_f\cup E_g),
        \end{align*}
        como $\mu(E_f\cup E_g)=0$, esto implica que $f+g\in M$. Ahora, si $x \in (-1,1)\setminus E_f$, $f(x)=f(-x)$, por tanto 
        \begin{align*}
            \alpha f(x)=\alpha f(-x) \text{ para todo } x \in (-1,1)\setminus E_f,
        \end{align*}
        y como $\mu(E_f)=0$, se tiene que $\alpha f\in M$. De esta manera $M$ es un subespacio de $L^2((-1,1))$.
        Ahora, recordemos que para $f,g \in L^2((-1,1))$, el producto interno $(f;g)$ está dado por 
        \begin{align*}
            (f,g)=\int_{-1}^1 f(x)g(x)\, dx.
        \end{align*}
        Como $M$ es un subespacio cerrado, para $f \in L^2((-1,1))$, $P_M f$ se caracteriza como la única $g\in M$ tal que $(f-g,h)=0$ para toda $h \in M$, es decir, para toda $h \in M$ se tiene
        \begin{align*}
            (f-g,h)=\int_{-1}^1 (f(x)-g(x))h(x)\, dx=0.
        \end{align*}
        Tomemos $g(x)=\dfrac{f(x)+f(-x)}{2}\in M$ para $x \in (-1,1)$, y veamos que\\ $g=P_M f$. Primero, es claro que $g \in M$, dado que 
        \begin{align*}
            g(-x)=\dfrac{f(-x)+f(-(-x))}{2}=\dfrac{f(x)+f(-x)}{2}=g(x),
        \end{align*}
        y para $h \in M$
        \begin{align*}
            (f-g,h)&=\int_{-1}^1\left(f(x)-\dfrac{f(x)+f(-x)}{2}\right)h(x)\, dx\\
            &=\int_{-1}^1\frac{f(x)-f(-x)}{2}h(x)\, dx\\
            &=\underbrace{\dfrac{1}{2}\int_{-1}^1f(x)h(x)\, dx}_{I_1}-\underbrace{\dfrac{1}{2}\int_{-1}^1 f(-x)h(x)\, dx}_{I_2}.
        \end{align*}
        Como $h\in M$, $h(x)=h(-x)$ para casi todo $x \in (-1,1)$, de esta manera, usando la sustitución $y=-x$ tenemos
        \begin{align*}
            I_2=\dfrac{1}{2}\int_{-1}^1 f(-x)h(x)\, dx=\dfrac{1}{2}\int_{-1}^1f(y)h(y)\, dy=I_1,
        \end{align*}
        de manera que 
        \begin{align*}
            (f-g,h)=I_1-I_2=I_1-I_1=0,
        \end{align*}
        por tanto, $(P_M f)(x)=g(x)=\dfrac{f(x)+f(-x)}{2}$.

        \item[(b)] Primero, note que podemos expresar $M$ usando el producto interno de $L^2((-1,1))$.
        \begin{align*}
            M=\left\{f \in L^2((-1,1)):\int_{-1}^1f(x)\, dx=0\right\}=\left\{f\in L^2((-1,1)): (f;1)=0\right\},
        \end{align*}
        donde $1$ denota la función constante $g(x)=1$. Antes de empezar a hacer los trámites, es importante mencionar que para $f \in L^2((-1,1))$, tiene sentido hallar la integral 
        \begin{align*}
            \int_{-1}^1f(x)\, dx,
        \end{align*}
        porque $\mu((-1,1))=2<\infty$, de donde se obtiene que $L^2((-1,1))\subset L^1((-1,1))$, más precisamente, usando la desigualdad de Cauchy-Schwarz, para \\
        $f \in L^2((-1,1))$ se tiene que 
        \begin{align*}
            \left|\int_{-1}^1f(x)\, dx\right|&=\left|\int_{-1}^1f(x)\cdot 1\, dx\right|\\
            &\leq \left(\int_{-1}^1|f(x)|^2\, dx\right)^{1/2}\left(\int_{-1}^1 1^2\, dx\right)^{1/2}\\
            &=\norm{f}_{L^2((-1,1))}(\mu((-1,1)))^{1/2}\\
            &=\sqrt{2}\norm{f}_{L^2((-1,1))}\\
            &<\infty.
        \end{align*}
        
        
        Note que $L^2((-1,1))\setminus M\neq \emptyset$, dado que para la función $g(x)=1$
        \begin{align*}
            (g,1)=\int_{-1}^1 g(x)\, dx=\int_{-1}^1 1 \, dx=2\neq 0,
        \end{align*}
        es decir, $g(x)=1\in L^2((-1,1))\setminus M$. Si $f \in L^2((-1,1))\setminus M$, entonces 
        \begin{align*}
            (f,1)\neq 0,
        \end{align*}
        definimos
        \begin{align*}
            \alpha=\left|\left(f,\dfrac{1}{\sqrt{2}}\right)\right|=\dfrac{1}{\sqrt{2}}|(f,1)|>0,
        \end{align*}
        y sea $B=B\left(f;\dfrac{\alpha}{2}\right)$, la bola en $L^2((-1,1))$ centrada en $f$ con radio $\dfrac{\alpha}{2}$. Sea \\
        $g \in B$, entonces $\norm{f-g}_{L^2((-1,1))}<\dfrac{\alpha}{2}$, además, por la desigualdad de Cauchy-Schwarz, tenemos que 
        \begin{align*}
            \left|\left(g-f,\dfrac{1}{\sqrt{2}}\right)\right|\leq \norm{g-f}_{L^2((-1,1))}\norm{\dfrac{1}{\sqrt{2}}}_{L^2((-1,1))}=\norm{g-f}_{L^2((-1,1))}<\dfrac{\alpha}{2},
        \end{align*}
        de esta manera
        \begin{align*}
            \left|\left(g,\dfrac{1}{\sqrt{2}}\right)\right|&=\left|\left(f,\dfrac{1}{\sqrt{2}}\right)+\left(g-f,\dfrac{1}{\sqrt{2}}\right)\right|\\
            &\geq \left|\left(f,\dfrac{1}{\sqrt{2}}\right)\right|-\left|\left(g-f,\dfrac{1}{\sqrt{2}}\right)\right|\\
            &>\alpha-\dfrac{\alpha}{2}\\
            &=\dfrac{\alpha}{2},
        \end{align*}
        de manera que 
        \begin{align*}
            \left|\left(g,\dfrac{1}{\sqrt{2}}\right)\right|>\dfrac{\alpha}{2}>0\\
            \Longrightarrow |(g,1)|>\dfrac{\sqrt{2}\alpha}{2}>0,
        \end{align*}
        luego $(g,1)\neq 0$ y así, $g \notin M$. De esta manera, tenemos que\\ $B\subset L^2((-1,1))\setminus M$, concluyendo que $M$ es cerrado. 

        Veamos ahora que $M$ es un subespacio de $L^2((-1,1))$. Sean $f,g \in M$ y $\alpha\in \mathbb{R}$, entonces
        \begin{align*}
            \int_{-1}^1 (f(x)+g(x))\, dx=\int_{-1}^1f(x)\, dx+\int_{-1}^1g(x)\, dx=0+0=0,
        \end{align*}
        es decir, $f+g\in M$, además
        \begin{align*}
            \int_{-1}^1\alpha f(x)\, dx=\alpha\int_{-1}^1 f(x)\, dx=\alpha\cdot 0=0,
        \end{align*}
        por tanto, $\alpha f\in M$. De esta manera, $M$ es un subespacio de $L^2((-1,1))$. 

        Hallemos ahora $P_M$. Sea $f\in L^2((-1,1))$, queremos encontrar $g\in M$ tal que $(f-g,h)=0$ para todo $h \in M$, es decir
        \begin{align*}
            \int_{-1}^1 ((f(x)-g(x))h(x)\, dx=0.
        \end{align*}
        Tomemos $g(x)=\displaystyle f(x)+C_f$ donde
        \begin{align*}
            C_f=-\dfrac{1}{2}\int_{-1}^1 f(y)\, dy,
        \end{align*}
        y veamos que $g=P_M f$. Primero veamos que, en efecto, $g \in M$ 
        \begin{align*}
            \int_{-1}^1g(x)\, dx&=\int_{-1}^1 (f(x)+C_f)\, dx\\
            &=\int_{-1}^1f(x)\, dx+C_f\int_{-1}^1\, dx\\
            &=\int_{-1}^1f(x)\, dx+2C_f\\
            &=\int_{-1}^1 f(x)\ dx-\int_{-1}^1f(y)\, dy\\
            &=0,
        \end{align*}
        por tanto $g \in M$. Ahora sea $h \in M$
        \begin{align*}
            (f-g,h)&=\int_{-1}^1 (f(x)-g(x))h(x)\, dx\\
            &=\int_{-1}^1 (f(x)-(f(x)+C_f))h(x)\, dx\\
            &=C_f\int_{-1}^1h(x)\, dx\\
            &=C_f\cdot 0\\
            &=0,
        \end{align*}
        de esta manera, $g=f+C_f=P_M f$.

        \item[(c)] Note que,
        \begin{align*}
            g\in M \Longleftrightarrow g(x)=0\, \text{ c.t.p } \, x\in (-1,0) \Longleftrightarrow \int_{-1}^0|g(x)|^2\, dx=0,
        \end{align*}
        además, 
        \begin{align*}
            \int_{-1}^0|g(x)|^2\, dx=\int_{-1}^1|g(x)|^2\chi_{(-1,0)}(x)\, dx=\int_{-1}^1|g(x)\chi_{(-1,0)}(x)|^2\, dx=\norm{\chi_{(-1,0)}g}_{L^2((-1,1))}^2.
        \end{align*}
        de manera que $g \in M$ si y sólo si $\norm{\chi_{(-1,0)}g}_{L^2((-1,1))}=0$. Sea $(f_n)$ una sucesión de funciones en $M$ tal que $f_n\to f$, Note que 
        \begin{align*}
            \norm{\chi_{(-1,0)} f_n-\chi_{(-1,0)}f}_{L^2(-1,1)}&=\norm{\chi_{(-1,0)}(f_n-f)}_{L^2((-1,1))}\\
            &=\left(\int_{-1}^1\left|\chi_{(-1,0)}(f_n(x)-f(x))\right|^2\right)^{1/2}\\
            &=\left(\int_{-1}^0 |f_n(x)-f(x)|^2\, dx\right)^{1/2}\\
            &\leq \left(\int_{-1}^1 |f_n(x)-f(x)|^2\, dx\right)^{1/2}\\
            &=\norm{f_n-f}_{L^2((-1,1))},
        \end{align*}
        de manera que $\chi_{(-1,0)}f_n\to \chi_{(-1,0)}f$ en $L^2((-1,1))$. Para ver que $f \in M$ verificamos que $\norm{\chi_{(-1,0)}f}_{L^2((-1,1))}=0$. Sea $\varepsilon>0$, existe $N_\varepsilon$ tal que si $n\geq N$, entonces 
        \begin{align*}
            \norm{\chi_{(-1,0)}f_n-\chi_{(-1,0)}f}_{L^2((-1,1))}=\norm{\chi_{(-1,0)}(f_n-f)}_{L^2((-1,1))}<\varepsilon,
        \end{align*}
        Para $n\geq N_\varepsilon$ tenemos
        \begin{align*}
            \norm{\chi_{(-1,0)}f}_{L^2((-1,1))}&=\norm{\chi_{(-1,0)}(f-f_n)+\chi_{(-1,0)}f_n}_{L^2((-1,1))}\\
            &\leq \norm{\chi_{(-1,0)}(f_n-f)}_{L^2((-1,1))}+\norm{\chi_{(-1,0)}f_n}_{L^2((-1,1))}\\
            &=\norm{\chi_{(-1,0)}(f_n-f)}_{L^2((-1,1))}\\
            &<\varepsilon,
        \end{align*}
        de esta manera $\norm{\chi_{(-1,0)}f}_{L^2((-1,1))}<\varepsilon$, y como $\varepsilon>0$ es arbitrario, obtenemos que 
        \begin{align*}
            \norm{\chi_{(-1,0)}f}_{L^2((-1,1))}=0,
        \end{align*}
        de donde obtenemos que $f(x)=0$ para casi todo $x \in (-1,0)$ y por tanto, $f\in M$. Así, concluimos que $M$ es cerrado.

        Veamos ahora que $M$ es un subespacio de $L^2((-1,1))$. Sean $f,g \in M$ y $\alpha\in \mathbb{R}$, sean 
        \begin{align*}
            E_f&=\{x\in (-1,0):f(x)\neq 0\}\\
            E_g&=\{x \in (-1,0):g(x)\neq 0\},
        \end{align*}
        como $f,g \in M$, $\mu(E_f)=\mu(E_g)=0$ y por tanto
        \begin{align*}
            0\leq \mu(E_f\cup E_g)\leq \mu(E_f)+\mu(E_g)=0+0=0,
        \end{align*}
        es decir, $\mu(E_f\cup E_g)=0$, entonces, para $x \in (-1,0)\setminus (E_f\cup E_g)$, $f(x)=0$ y $g(x)=0$, por tanto
        \begin{align*}
            f(x)+g(x)=0 \text{ para todo } x \in (-1,0)\setminus (E_f\cup E_g),
        \end{align*}
        y como $\mu(E_f\cup E_g)=0$, tenemos que $f+g\in M$. Si $x \in (-1,0)\setminus E_f$, entonces $f(x)=0$, por tanto
        \begin{align*}
            \alpha f(x)=\alpha\cdot 0=0 \text{ para todo } x \in (-1,0)\setminus E_f,
        \end{align*}
        nuevamente, como $\mu(E_f)=0$, se tiene que $f \in M$. De esta manera, concluimos que $M$ es un subespacio de $L^2((-1,1))$. 

        Calculemos la proyección ortogonal $P_M$. Sea $f\in L^2((-1,1))$, queremos encontrar $g \in M$ tal que $(f-g,h)=0$ para toda $h \in M$. Tomemos $g=\chi_{[0,1)}f$ y veamos que $g=P_M f$. Claramente $g(x)=0$ para todo $x \in (-1,0)$, por lo que $g \in M$. Sea $h \in M$, entonces
        \begin{align*}
            (f-g,h)&=\int_{-1}^1(f(x)-\chi_{[0,1)}(x)f(x))h(x)\, dx\\
            &=\int_{-1}^1\chi_{(-1,0)}(x)f(x)h(x)\, dx\\
            &=\int_{-1}^0f(x)h(x)\, dx\\
            &=0,
        \end{align*}
        dado que $h(x)=0$ para casi todo $x\in (-1,0)$. Así, concluimos que\\ $g=\chi_{[0,1)}f=P_Mf$.
        \end{enumerate}

        \item[(II)] \begin{enumerate}
            \item[(a)] Veamos que $K$ es cerrado en $L^2(\Omega)$. Primero, note que
        \begin{align*}
            K=\left\{f \in L^2(\Omega):(f;1)\geq 1\right\}.
        \end{align*}
        Como $\Omega$ es un abierto acotado, tenemos que $0<\mu(\Omega)<\infty$, además
        \begin{align*}
            \norm{1}_{L^2(\Omega)}=\left(\int_{\Omega}|1|^2\, dx\right)^{1/2}=(\mu(\Omega))^{1/2}.
        \end{align*}
        
        Sea $g\in L^2(\Omega)\setminus K$, es decir, $\alpha:=(g,1)<1$. Tomemos $\varepsilon>0$ tal que $\alpha+\varepsilon<1$ y  $\delta=\dfrac{\varepsilon}{(\mu(\Omega))^{1/2}}$. Consideremos $B=B(g,\delta)$, la bola en $L^2(\Omega)$ centrada en $g$ y de radio $\delta$ y sea $f \in B$, entonces $\norm{f-g}_{L^2(\Omega)}<\delta$. Por la desigualdad de Cauchy-Schwarz y lo anterior, se tiene que 
        \begin{align*}
            (f,1)&=(f-g,1)+(g,1)\\
            &\leq \norm{f-g}_{L^2(\Omega)}\norm{1}_{L^2(\Omega)}+\alpha\\
            &<\delta (\mu(\Omega))^{1/2}+\alpha\\
            &=\dfrac{\varepsilon}{(\mu(\Omega))^{1/2}}(\mu(\Omega))^{1/2}+\alpha\\
            &=\varepsilon+\alpha\\
            &<1,
        \end{align*}
        de manera que $(f;1)<1$ y así, $B\subset L^2(\Omega)\setminus K$. De esta manera, concluimos que $K$ es cerrado en $L^2(\Omega)$.

        Veamos que $K$ es convexo. Sean $f,g \in K$ y $t \in [0,1]$, entonces
        \begin{align*}
            \int_{\Omega}(tf(x)+(1-t)g(x))\, dx=t\int_{\Omega}f(x)\, dx+(1-t)\int_{\Omega}g(x)\, dx\geq t+(1-t)=1,
        \end{align*}
        luego $tf+(1-t)g \in K$, concluyendo que $K$ es convexo. 

        \item[(b)] Procedemos a calcular $P_K$. Sea $f \in L^2(\Omega)$, queremos encontrar $g\in K$ tal que $(f-g;h-g)\leq 0$ para toda $h \in K$. Proponemos 
        \begin{align*}
            g(x)=f(x)+\chi_{(-\infty,1)}(C_f)\dfrac{1-C_f}{\mu(\Omega)},
        \end{align*}
        donde 
        \begin{align*}
            C_f=\int_\Omega f(y)\, dy.
        \end{align*}
        Veamos que, en efecto, $g=P_K f$. Tenemos dos casos
        \begin{itemize}
            \item Si $f \in K$, tenemos que $C_f\geq 1$ y por tanto $\chi_{(-\infty,1)}(C_f)=0$, por tanto, $g=f$ y así, $g \in K$ y $(f-g,h-g)=(0,h-f)=0\leq 0$, es decir, $g=f=P_Kf$.

            \item Si $f\notin K$, tenemos que $C_f<1$, por tanto $\chi_{(-\infty,1)}(C_f)=1$. Veamos que $g \in K$
            \begin{align*}
                \int_{\Omega}g(x)\, dx&=\int_{\Omega}\left(f(x)+\chi_{(-\infty,1)}(C_f)\dfrac{1-C_f}{\mu(\Omega)}\right)\, dx\\
                &=\int_{\Omega}\left(f(x)+\dfrac{1-C_f}{\mu(\Omega)}\right)\, dx\\
                &=\int_{\Omega}f(x)\, dx+\dfrac{1-C_f}{\mu(\Omega)}\int_{\Omega}\, dx\\
                &=C_f+(1-C_f)\\
                &=1.
            \end{align*}
            por tanto, $g\in K$.
            
            Note que, como $C_f<1$, entonces $1-C_f>0$ y, por tanto, $\dfrac{1-C_f}{\mu(\Omega)}>0$. Sea $h \in K$ cualquiera, entonces $C_h\geq 1$ y por tanto, $C_h-1\geq 0$
            \begin{align*}
                (f-g,h-g)&=\left(-\dfrac{1-C_f}{\mu(\Omega)},h-f-\dfrac{1-C_f}{\mu(\Omega)}\right)\\
                &=-\dfrac{1-C_f}{\mu(\Omega)}\left(1,h-f-\dfrac{1-C_f}{\mu(\Omega)}\right)\\
                &=-\dfrac{1-C_f}{\mu(\Omega)}\int_{\Omega}\left(h(x)-f(x)-\dfrac{1-C_f}{\mu(\Omega)}\right)\, dx\\
                &=-\dfrac{1-C_f}{\mu(\Omega)}\left(\int_{\Omega}h(x)\, dx-\int_{\Omega}f(x)\, dx-\dfrac{1-C_f}{\mu(\Omega)}\int_{\Omega}\, dx\right)\\
                &=-\dfrac{1-C_f}{\mu(\Omega)}(C_h-C_f-(1-C_f))\\
                &=-\dfrac{1-C_f}{\mu(\Omega)}(C_h-1)\\
                &\leq 0.
            \end{align*}
            de esta manera $g=P_K f$.
        \end{itemize}
        Así, podemos concluir que 
        \begin{align*}
            (P_K f)(x)=f(x)+\chi_{(-\infty,1)}(C_f)\dfrac{1-C_f}{\mu(\Omega)}, \hspace{5mm} C_f=\int_{\Omega}f(y)\, dy.
        \end{align*}
        \end{enumerate}
    \end{enumerate}
    \end{proof}


\textbf{Ejercicio 14.} Sea $H$ un espacio de Hilbert y $A\in L(H)=L(H,H)$ (el conjunto de funciones lineales continuas de $H$ en $H$).
\begin{enumerate}
    \item[(I)] Para $y \in H$ fijo, muestre que el funcional $\Phi_y:H\to \mathbb{R}$ dado por $(Ax,y)$ es lineal y continuo. Deduzca que existe un único elemento en $H$ que denotaremos por $A^\star y$, tal que 
    \begin{align*}
        (Ax,y)=(x,A^\star y) \text{ para todo } x \in H.
    \end{align*}
    \item[(II)] Muestre que $A^\star \in L(H)$. $A^\star$ se llama el adjunto de $A$.
    \item[(II)] Verifique que $(A^\star)^\star$ y $\norm{A^\star}=\norm{A}$.
\end{enumerate}
\begin{proof}
    \begin{enumerate}
        \item [(I)]Veamos que $\Phi_y$ es lineal. Como H es de Hilbert, y usando que $A \in \mathcal{L}(H)$ se sigue que para todo $x,z \in H$ y $\lambda \in \mathbb{R}$:
\begin{align*}
\Phi_y(x+\lambda z) &= ((Ax+\lambda z), y) \\
&= (Ax+\lambda Az, y) \\
&= (Ax,y) + (\lambda Az,y) \\
&= (Ax,y) + \lambda (Az,y) \\
&= \Phi_y(x) + \lambda \Phi_y(z)
\end{align*}

Para ver que $\Phi_y$ es continuo, como $A \in \mathcal{L}(H)$, existe una constante $M\geq 0$ tal que $\norm{Ax}\leq M\norm{x}$ para todo $x \in H$. Así, dado $x \in H$
\begin{align*}
|\Phi_y x| &= |(Ax, y)| \\
&\le \|Ax\| \|y\| \\
&\le M \|x\| \|y\|.
\end{align*}
con lo que se concluye que $\Phi_y \in H^{\star}$.

Ahora, por el teorema de representación de Riesz, existe un único elemento en $H$, llámelo $A^{\star}y$, tal que
\begin{align*}
(Ax,y) &= \langle \Phi_y, x \rangle \\
&= (A^{\star}y, x) \\
&= (x, A^{\star}y) \quad \forall x \in H.
\end{align*}
en donde en la última igualdad se usa que $(\cdot, \cdot)$ es simétrico.
\item[(II)] Primero veamos que $A^\star$ es lineal. Para ello, dados $y, z \in H$ y $\lambda \in \mathbb{R}$,
queremos ver que
\[
(x, A^\star(\lambda y+z)) = (x, \lambda A^\star y + A^\star z) \quad \forall x \in H
\]
Así, dado $x \in H$
\begin{align*}
(x, A^\star(\lambda y+z)) &= (Ax, \lambda y+z) \\
&= \lambda (Ax,y) + (Ax,z) \\
&= \lambda (x, A^\star y) + (x, A^\star z) \\
&= (x, \lambda A^\star y) + (x, A^\star z) \\
&= (x, \lambda A^\star y + A^\star z)
\end{align*}
Por tanto, como
\[
(x, A^\star(\lambda y+z)) = (x, \lambda A^\star y + A^\star z) \quad \forall x \in H
\]
entonces $A^\star(\lambda y+z) = \lambda A^\star y + A^\star z$.

Veamos que $A^\star$ es continuo.
Para esto, usaremos el Teorema del Gráfico cerrado. Con el fin de evitar confusiones en la notación, usaremos $(\cdot,\cdot)_H$ para denotar el producto interior de H, mientras que la notación de pareja ordenada $(\cdot,\cdot) \in H^2$ se mantiene igual.

Sea $(x,y) \in \overline{\mathcal{G}(A^\star)}$. Luego existe una sucesión $\{ (x_n, A^\star x_n) \}$ en $\mathcal{G}(A^\star)$ tal que
\begin{align*}
x_n &\xrightarrow{n \to \infty} x \\
A^\star x_n &\xrightarrow{n \to \infty} y
\end{align*}
Sea $z \in H$, luego
\[
(Az,x_n)_H = (z, A^\star x_n)_H
\]
Ahora, notemos que $(\cdot,\cdot)_H: H^2 \longrightarrow \mathbb{R}$ es continuo; pues dado $\{u_n\}, \{v_n\}$ sucesiones en H tales que
\begin{align*}
u_n &\longrightarrow u \\
v_n &\longrightarrow v
\end{align*}
se tiene que
\begin{align*}
|(u_n, v_n)_H - (u, v)_H| &= |(u_n, v_n)_H - (u, v_n)_H + (u, v_n)_H - (u, v)_H| \\
&\le |(u_n - u, v_n)_H| + |(u, v_n - v)_H| \\
&\le \|u_n - u\| \|v_n\| + \|u\| \|v_n - v\|
\end{align*}
Y como dado $\varepsilon>0$, existe $N>0$ tal que $\|u_n - u\| < \varepsilon$, $\|v_n - v\| < \varepsilon$, siempre que $n\geq N$.
Además, como $\{v_n\}$ es acotada por ser convergente, existe $M>0$ tal que $\|v_n\| \le M \quad \forall n \ge 1$. Entonces
\[
|(u_n, v_n)_H - (u, v)_H| < M\varepsilon + \|u\|\varepsilon
\]
Y por tanto, $(\cdot,\cdot)_H$ es continuo.
Usando esto y que
\[
(Az, x_n)_H = (z, A^\star x_n)_H
\]
haciendo tender $n \to \infty$
\begin{align*}
(Az, x)_H &= (z, y)_H \\
(z, A^\star x)_H &= (z, y)_H \\
(z, A^\star x - y) &= 0
\end{align*}
Como esto ocurre para cualquier $z \in H$,
\[
A^\star x = y
\]
de modo que $(x,y)=(x, A^\star x) \in \mathcal{G}(A^\star)$ y por tanto, $\mathcal{G}(A^\star)$ es cerrado, con lo cual concluimos que $A^\star$ es continuo y así, $A^\star \in \mathcal{L}(H)$.
\item[(III)] Sea $y \in H$. Así:
\begin{align*}
(((A^\star)^\star - A)y, x) &= ((A^\star)^\star y - Ay, x) \\
&= ((A^\star)^\star y, x) - (Ay, x) \\
&= (x, (A^\star)^\star y) - (Ay, x) \\
&= (A^{\star}x, y) - (y, A^{\star}x) \\
&= (A^\star x, y) - (A^{\star}x, y)=0
\end{align*}
para todo $x \in H$. Por tanto, $(A^\star)^\star y = Ay$. Pero como $y \in H$ es arbitrario, $\displaystyle(A^\star)^\star = A$.

Finalmente, mostremos que $\|A^\star\| \le \|A\|$ y $\|A\| \le \|A^\star\|$. \\

\checkmark$\|A^\star\| \le \|A\|$: Sea $x \in H$, con $\|x\|=1$. Así
\begin{align*}
(A^\star x, A^\star x) &= (A A^\star x, x) \\
&\le \|A A^\star x\| \|x\| \\
&\le \|A\| \|A^\star x\| 
\end{align*}
Como $(A^\star x, A^\star x) = \|A^\star x\|^2$, se sigue que
\[
\|A^\star x\| \leq \|A\|
\]
pues en el caso que $\|A^\star x\|=0$, la desigualdad se tiene trivialmente. De esta forma
\[
\|A^\star\| = \sup_{\substack{x \in H \\ \|x\|=1}} \|A^\star x\| \le \|A\|
\]

\checkmark$\|A\| \le \|A^\star\|$: 
Por lo hecho en la primera parte, $A = (A^\star)^\star$.
Así (y usando lo probado anteriormente)
\[
\|A\| = \|(A^\star)^\star\| \le \|A^\star\|. 
\] 
Por lo tanto, $\norm{A}=\norm{A^{\star}}$
 
    \end{enumerate}
\end{proof}
\textbf{Ejercicio 15.} Sea $H$ un espacio de Hilbert y $M\subseteq H$ un subespacio cerrado. Considere la proyección ortogonal $P_M$. Muestre que 
\begin{enumerate}
    \item[(I)] $P_M$ es lineal.
    \item[(II)] $P_M^2=P_M$.
    \item[(III)] $P_M^\star=P_M$, donde $P_M^\star$ denota el operador adjunto de $P_M$.
    \item[(IV)] $\text{Rango}(P_M)=M$ y $\text{Kernel}(P_M)=M^\perp$.
    \item[(V)] Suponga que $P \in L(H)$. Entonces $P$ es la proyección sobre un subespacio cerrado si y sólo si $P=P^2=P^\star$.
\end{enumerate}
\begin{proof} 
  \begin{enumerate}
    \item[(I)] Veamos que $P_{M}$ es lineal, note que como $M$ es un subespacio cerrado, sabemos que dado un $f\in H$ existe un único $u\in M$ tal que
      \begin{align*}
        |f-u|=\min_{v\in M}|f-v|=dist(f,M).
      \end{align*}
      más aún
      \begin{align*}
        (f-u,v)&= 0&&\text{para todo $v\in M$.}
      \end{align*}
      Siendo así, suponga $f,g\in H$ y $\lambda$ escalar, entonces sabemos que existen $u_{1},u_{2}\in M$ tales que  
      \begin{align*}
        (f-u_{1},v)&=0,\\
        (g-u_{2},v)&=0&&\text{para todo $v\in M$.}
      \end{align*}
      Luego se puede inferir usando la linealidad del producto interno que
      \begin{align*}
        0&=(f-u_{1},v)+\lambda(g-u_{2},v)=(f+\lambda g-(u_{1}+\lambda u_{2}),v) &&\text{para todo $v\in M$.}
      \end{align*}
      De lo que se concluye que si $P_{M}(f)=u_{1}, P_{M}(g)=u_{2}$, entonces $P_{M}(f+\lambda g)&=P_{M}(f)+\lambda P_{M}(g)$, es decir, el operador proyección ortogonal $P_{M}$ es un operador lineal. 
    \item[(II)] Note que dado $f\in H$, entonces existe un único $u\in M$ tal que $P_{M}(f)=u\in M$, además, como $u\in M$, entonces $P_{M}(u)=u$, ya que
    \begin{align*}
      (u-u,v)=(0,v)&=0&&\text{para todo $v\in M$.}
    \end{align*}
    De lo que se concluye que
    \begin{align*}
      P_{M}^{2}(f)&=P_{M}\left( P_{M}(f) \right),\\
      &=P_{M}(u),\\
      &=u,\\
      &=P_{M}(f).
    \end{align*}
    Lo que concluye el resultado.
    \item[(III)] como $H$ es un espacio de Hilbert y $P_{M}:H\to H$, entonces por el teorema de representación de Riesz sabemos que el operador $P_{M}^{\star}$ cumple que dados $x,y\in H$ se satisface que 
    \begin{align*}
      \left(P_{M}(x),y\right)&=\left(x,P_{M}^{\star}(y)\right)
    \end{align*}
    \item[(IV)]
  \end{enumerate}
\end{proof}


\section{Operadores Compactos y Teorema Espectral}

\textbf{Ejercicio 3.} Considere los operadores de desplazamiento $S_r,S_l\in L(\ell^2)$, donde si\\ $x=(x_1,x_2,x_3,...)\in \ell^2$, estos se definen como
\begin{align*}
    S_rx:=(0,x_1,x_2,x_3,...),
\end{align*}
y
\begin{align*}
    S_lx=(x_2,x_3,x_4,...).
\end{align*}
$S_r$ se conoce como el desplazamiento a derecha y $S_l$ como el desplazamiento a izquierda.

\begin{enumerate}
    \item[(a)] Determinar las normas $\norm{S_r}$ y $\norm{S_l}$.
    \item[(b)] Muestre que $EV(S_r)=\emptyset$.
    \item[(c)] Muestre que $\sigma(S_r)=[-1,1]$.
    \item[(d)] Muestre que $EV(S_l)=(-1,1)$. Encuentre el subespacio propio correspondiente.
    \item[(e)] Muestre que $\sigma(S_l)=[-1,1]$.
    \item[(f)] Determine los adjuntos $S_r^\star$ y $S_l^\star$.
\end{enumerate}

\textbf{Ejercicio 4.} Sea $1\leq p<\infty$ y consideremos el espacio $L^p((0,1))$. Dado $u \in L^p((0,1))$, definimos
\begin{align*}
    Tu(x)=\int_{0}^xu(t)\, dt.
\end{align*}
\begin{enumerate}
    \item[(a)] Demuestre que $T \in \mathcal{K}(L^p((0,1)))$. 
    \item[(b)] Determine $EV(T)$ y $\sigma(T)$.
    \item[(c)] Dé una fórmula explícita para $(T-\lambda I)^{-1}$ cuando $\lambda\in \rho(T)$.
    \item[(d)] Determine $T^\star$.
\end{enumerate}

\begin{proof}
    \begin{enumerate}
        \item[(a)] Veamos primero que $T\in \mathcal{L}(L^p((0,1)))$. Por simplicidad, denotaremos $\norm{\cdot}_{L^p((0,1))}=\norm{\cdot}_p$. Es claro que $T$ es un operador lineal, además, dada $u \in L^p((0,1))$
        \begin{align*}
            \norm{Tu}_p&=\left(\int_{0}^1 |Tu(x)|^p\, dx\right)^{1/p}\\
            &=\left(\int_0^1 \left|\int_0^x u(t)\, dt\right|^p\, dx\right)^{1/p}\\
            &\leq \left(\int_{0}^1\left(\int_0^1 |u(t)|\, dt\right)^p\, dx\right)^{1/p}\\
            &=\int_0^1 |u(t)|\, dt\\
            &=\norm{u}_1\\
            &\leq \norm{u}_p.
        \end{align*}
        de manera que $T$ es acotado. Veamos ahora que $T \in \mathcal{K}(L^p((0,1)))$. Sea
        \begin{align*}
            B=\left\{f \in L^p((0,1)): \norm{f}_p\leq 1\right\}.
        \end{align*}
        Si $u \in B$, por lo hecho anteriormente, tenemos que $\norm{Tu}_p\leq \norm{u}_p\leq 1$, es decir, $T(B)\subset B$.

        \item[(b)] Sea $\lambda\in \mathbb{R}$, si $u\in L^p((0,1))$ es diferenciable, tenemos que 

        \item[(c)]

        \item[(d)] Por definición, tenemos
        \begin{align*}
            T^\star: (L^p((0,1)))^\star &\longrightarrow (L^p((0,1)))^\star\\
            \xi &\longmapsto T^\star\xi: L^p((0,1)) \longrightarrow \mathbb{R}\\
            & \hspace{31.5mm} f \longmapsto \langle T^\star\xi; f\rangle:=\langle \xi;Tf\rangle. 
        \end{align*}
        Por el Teorema de Representación de Riesz, dado $\xi \in (L^p((0,1)))^\star$, existe una única $g_\xi \in L^{p'}((0,1))$ tal que 
        \begin{align*}
            \langle \xi;h\rangle=\int_0^1 g_\xi(x)h(x)\, dx,
        \end{align*}
        para toda $h \in L^p((0,1))$, donde $\dfrac{1}{p}+\dfrac{1}{p'}=1$. De manera Análoga, como\\ $T^\star \xi \in (L^p((0,1)))^\star$, existe una única $G_\xi\in L^{p'}((0,1))$, cumpliendo la misma relación con $p$, tal que 
        \begin{align*}
            \langle T^\star \xi;h\rangle=\int_0^1 G_\xi(x)h(x)\, dx,
        \end{align*}
    \end{enumerate}
    De esta manera, por la definición del operador adjunto $T^\star$, dada $f \in L^p((0,1))$ se tiene
    \begin{align*}
        \langle T^\star \xi;f\rangle=\int_{0}^1G_\xi(x)f(x)\, dx=\int_0^1 g_\xi(x)Tf(x)\, dx=\langle \xi;Tf\rangle,
    \end{align*}
    por definición y usando el Teorema de Fubini, dado que $f$ es arbitraria, obtenemos
    \begin{align*}
        \int_0^1 g_\xi(x)Tf(x)\, dx&=\int_0^1 g_\xi(x)\int_0^x f(t)\, dt\, dx\\
        &=\int_0^1 \int_0^x g_\xi(x)f(t)\, dt\, dx\\
        &=\int_0^1 \int_t^1 g_\xi(x)f(t)\, dx\, dt\\
        &=\int_0^1f(t)\int_{t}^1g_\xi(x)\, dx\, dt.
    \end{align*}
    de manera que
    \begin{align*}
        G_\xi(t)=\int_t^1 g_\xi(x)\, dx=\int_{0}^1 g_{\xi}(x)\chi_{(t,1)}(x)\, dx=\langle \xi;\chi_{(t,1)}\rangle,
    \end{align*}
    así, 
    \begin{align*}
        \langle T^\star \xi;f\rangle=\int_0^1 \langle \xi;\chi_{(x,1)}\rangle f(x)\, dx
    \end{align*}
\end{proof}

\textbf{Ejercicio 6.} Considere $g\in L^\infty(\mathbb{R})\cap C(\mathbb{R})$ (es decir, $g$ es continua y acotada). Definimos el operador de multiplicación $M_g:L^2(\mathbb{R})\to L^2(\mathbb{R})$ dado por
\begin{align*}
    M_g(f)(x)=g(x)f(x).
\end{align*}
\begin{enumerate}
    \item[(a)] Muestre que $\sigma(M_g)=\overline{\{g(x):x\in \mathbb{R}\}}$.
    \item[(b)] ¿Es el operador $M_g$ compacto?
\end{enumerate}
